\documentclass[a4paper]{article}

\usepackage{graphicx}
\usepackage{amsmath,amssymb}
\usepackage{minted}
\usepackage{multirow}
\usepackage{float}
\usepackage{todonotes}
\usepackage{forest}
\usepackage{xstring}
\usepackage{xcolor}
\usepackage[most]{tcolorbox}
\usepackage{xparse}
\usepackage{tikz}

\usetikzlibrary{graphs,graphdrawing}
\usegdlibrary{layered}

\input{/home/donkarlo/Dropbox/repo/nd_ai_project/src/nd_ai/knoweldge_representation/language/formal/computer/programming/tex/latex/code_clips/external_graphviz.tex}
\input{/home/donkarlo/Dropbox/repo/nd_ai_project/src/nd_ai/knoweldge_representation/language/formal/computer/programming/tex/latex/parts/control_sequence/kind/semtag/semtag.tex}
\input{/home/donkarlo/Dropbox/repo/nd_ai_project/src/nd_ai/knoweldge_representation/language/formal/computer/programming/tex/latex/code_clips/seehere.tex}

\title{Self}
\date{}

\begin{document}
	\maketitle
	
	\section{Neural basis of self}
		\seehere{https://en.wikipedia.org/wiki/Neural_basis_of_self}
	
	\section{From autobiographical memory}
		In \cite{brewer-1986-what-is-autobiographical-memory} page 26, the parts of self is mentioned and is summarized as
		The working self, often referred to as just the 'self', is a set of active personal goals or self-images organized into goal hierarchies. These personal goals and self-images work together to modify cognition and the resulting behavior so an individual can operate effectively in the world.
		\seehere{https://en.wikipedia.org/wiki/Autobiographical_memory?utm_source=chatgpt.com\#Working_self}
	
	\cite{brewer-1986-what-is-autobiographical-memory} includes the self in the following hierarchy:
	
	\begin{figure}[H]
		\centering
		\begin{forest}
			for tree={
				grow=south,
				align=center
			}
			[Individual
				[Self
					[Ego\\(experience stream)]
					[Self-schema\\(knowledge about self)]
					[Autobiographical memory\\(episodes + facts)]
				]
			]
		\end{forest}
	\end{figure}
	
	\begin{itemize}
		\item \textbf{Autobiographical memory}
		\begin{itemize}
			\item Defined as memory for information related to the self.
		\end{itemize}
		
		\item \textbf{Structure of the self}
		\begin{itemize}
			\item The self consists of:
			\begin{itemize}
				\item an experiencing ego,
				\item a self-schema,
				\item and associated personal memories and autobiographical facts.
			\end{itemize}
		\end{itemize}
		
		\item \textbf{Ego}
		\begin{itemize}
			\item The conscious experiencing entity.
			\item Represents moment-to-moment subjective experience.
			\item Personal memories correspond to the ego’s experiences over time.
		\end{itemize}
		
		\item \textbf{Self-schema}
		\begin{itemize}
			\item A cognitive structure containing general knowledge about the self.
			\item Organized, relatively stable, and partly unconscious.
			\item Integrates private and publicly observable information about the self.
		\end{itemize}
		
		\item \textbf{Self}
		\begin{itemize}
			\item A complex structure including ego, self-schema, and self-related long-term memories.
		\end{itemize}
		
		\item \textbf{Individual}
		\begin{itemize}
			\item A broader entity including the self, body, depersonalized knowledge, and various cognitive and motor skills.
		\end{itemize}
	\end{itemize}

	\bibliographystyle{plain}
	\bibliography{/home/donkarlo/Dropbox/repo/nd_shared_research/bibliography/bibliography.bib}

\end{document}
